\begin{trapbox}
	\begin{description}[style=nextline, leftmargin=2.8em, labelsep=0.5em, itemsep=0.35em]
		\item[\textbf{\textcolor{CTrap}{易错点一（公比条件忽略）}}] 使用公式前未检查 $q = 1$，直接代入导致分母为零。
		\item[\textbf{\textcolor{CTrap}{正确理解（公比必须非一）：}}] 公式仅适用于 $q \neq 1$；当 $q = 1$ 时，$b_n$ 为常数列，求和为等差数列前 $n$ 项和乘以 $b_1$。
		\item[\textbf{\textcolor{CTrap}{错因分析（盲目套用）：}}] 记忆公式时忽略了前提条件，将适用范围推广到所有等比情况。
	\end{description}

	\medskip
	\begin{description}[style=nextline, leftmargin=2.8em, labelsep=0.5em, itemsep=0.35em]
		\item[\textbf{\textcolor{CTrap}{易错点二（指数错误）}}] 公式第二项中误写为 $(1-q^{n})$ 或 $(1-q^{n-2})$，导致求和结果偏差。
		\item[\textbf{\textcolor{CTrap}{正确理解（指数核对）：}}] 第二项分子应为 $d\,b_1 q\,(1-q^{\,n-1})$，其中指数为 $n-1$，源于错位相减后等比部分有 $n-1$ 项。
		\item[\textbf{\textcolor{CTrap}{错因分析（推导不清）：}}] 对错位相减过程不熟悉，误以为等比部分项数为 $n$ 或 $n-2$。
	\end{description}

	\medskip
	\begin{description}[style=nextline, leftmargin=2.8em, labelsep=0.5em, itemsep=0.35em]
		\item[\textbf{\textcolor{CTrap}{易错点三（符号处理）}}] 当 $q$ 为负数时，代入公式时未加括号，导致指数运算正负出错。
		\item[\textbf{\textcolor{CTrap}{正确理解（括号保护）：}}] 将 $q$ 的值（含负号）整体代入公式，例如 $q = -2$ 时，$q^{\,n}$ 须写为 $(-2)^{n}$，不可写成 $-2^{n}$。
		\item[\textbf{\textcolor{CTrap}{错因分析（运算习惯）：}}] 平时处理正数习惯省略括号，忽略负数的奇偶次幂差异。
	\end{description}
\end{trapbox}
